% web: math. Most fixes mark something tex4ht would lose, for the Lua
% passes (lua/web/) to finish in the HTML. Redefined commands are
% \protected, so that they survive moving arguments (section titles in the
% table of contents, captions).
\makeatletter

% --- delimiters -------------------------------------------------------------

% tex4ht wraps \bigl( .. \bigr) in an <mrow>, which breaks when the pair spans
% a line break in align; emit just the sized <mo>, with min/maxsize (mathsize
% would stretch the normal glyph, thickening its stroke). TeX's \big and co.
% pick delimiters 1.2em, 1.8em, 2.4em and 3em tall, in ems of the extension
% font (delimiters.extension_scale scales them to it); ask for a hair under,
% since MathJax's 1.2em glyphs fall just short of 1.2em and it would skip to
% the next size (as its own \big does). Not fences: MathJax treats a row
% between stretchy fences as \left..\right, which gets space around it
\newcommand\web@sizeddelim[5]{% big; l or r; mathopen or mathclose; open or close; height
    \Configure{#1#2}
        {\bgroup\Configure{#1}{}{}{}{}%
         \Configure{#3}{*}
            {<mo class="MathClass-#4" fence="false" stretchy="true" symmetric="true" minsize="#5" maxsize="#5">}{</mo>}
            {\Configure{#3}{}{}{}{}}}
        {\egroup}%
}
\newcommand\web@bigdelim[2]{% big, Big, ..; height
    \web@sizeddelim{#1}{l}{mathopen}{open}{#2}%
    \web@sizeddelim{#1}{r}{mathclose}{close}{#2}%
    % \bigm| and co. are relations, as in TeX. tex4ht's own (a configured
    % \mathrel) loses its size inside another \mathrel (\mathrel{\Bigm|}),
    % so mark the delimiter for delimiters.middle to size
    \protected\expandafter\def\csname #1m\endcsname##1{%
        \HCode{<mrow class="bigm" data-mjx-texclass="REL" data-size="#2">}##1\HCode{</mrow>}}%
}
\web@bigdelim{big}{1.19em}
\web@bigdelim{Big}{1.79em}
\web@bigdelim{bigg}{2.39em}
\web@bigdelim{Bigg}{2.99em}

% \WebPairedDelimiters{abs,norm}: mathtools' \DeclarePairedDelimiter wraps
% \abs[\big]{x}'s delimiters in \mathopen{..} and \mathclose{..}, inside
% which tex4ht drops their size, and \abs*{x}'s in \mathopen{}\mathclose
% \bgroup .. \aftergroup\egroup, which it can't follow with mleftright; leave
% the delimiters of the commands listed bare. (TeX can't list the commands
% a document declared, so a project config names them)
\ExplSyntaxOn
\NewDocumentCommand \WebPairedDelimiters { m }
  {
    \cs_if_exist:NTF \reDeclarePairedDelimiterInnerWrapper
      {
        \clist_map_inline:nn {#1}
          {
            \exp_args:Nc \reDeclarePairedDelimiterInnerWrapper {##1} { nostarscaled } { ####1 ####2 ####3 }
            \exp_args:Nc \reDeclarePairedDelimiterInnerWrapper {##1} { star } { ####1 ####2 ####3 }
          }
      }
      { \PackageWarning { web } { \token_to_str:N \WebPairedDelimiters \c_space_tl needs~mathtools } }
  }
\ExplSyntaxOff

% mleftright makes \left/\right \mathopen{}\mathclose\bgroup\left ..
% \right\aftergroup\egroup, which tex4ht can't follow; it drops the MathML
% for everything inside \left..\right, bmatrix and cases. tex4ht's own
% \left/\right call the definitions from before \begin{document}, so point
% them back at the originals
\@ifpackageloaded{mleftright}{%
    \expandafter\let\csname o:left:\endcsname\mleftright@OrgLeft
    \expandafter\let\csname o:right:\endcsname\mleftright@OrgRight
}{}

% --- spaces -----------------------------------------------------------------

% tex4ht writes \  (an ordinary, breakable space in TeX) as a no-break space,
% like ~; write a space (in math, tex4ht's own space)
\Configure{ }
    {\ifmathml \ifmtext \space \else
        \HCode{<mspace width="0.33em" class="nbsp"></mspace>}\fi
     \else \HCode{ }\fi}
% tex4ht finds spaces from the gaps between glyphs, and misses the one after
% inline math that ends in a symbol (\}, \infty, \mathcal{H}); write out the
% space that follows inline math (the browser collapses a doubled one)
\Configure{$}
    {\Configure{@math}{display="inline"}\DviMath}
    {\EndDviMath\aftergroup\web@mathspace}
    {}
\def\web@mathspace{\futurelet\web@next\web@mathspace@check}
\def\web@mathspace@check{\ifx\web@next\@sptoken\HCode{ }\fi}

\@ifpackageloaded{amsmath}{%
    % amsmath's \colon: a little space, an ordinary colon, then more space
    % (tex4ht writes a plain colon, which TeX spaces as a relation)
    \protected\def\colon{\HCode{<mspace width="0.111em"></mspace><mo class="MathClass-ord" data-mjx-texclass="ORD" stretchy="false">:</mo><mspace width="0.333em"></mspace>}}%
    % \pmod sets "mod" in the operator font, which tex4ht doesn't see, so it
    % comes out italic, and drops the \mkern after it
    \protected\def\pmod#1{\pod{\mathrm{mod}\HCode{<mspace width="0.333em"></mspace>}#1}}%
}{}

% relations tex4ht writes as a look-alike: \parallel as \| (which is
% ordinary in TeX; classes.tokens tells them apart by this mark), and
% amssymb's slanted \leqslant and \geqslant as ≤ and ≥
\newcommand\web@relation[2]{% command, code point
    \protected\def#1{\mathrel{\HCode{<mo class="MathClass-rel" data-mjx-texclass="REL">\string&\#x#2;</mo>}}}%
}
\web@relation\parallel{2225}
\@ifpackageloaded{amssymb}{%
    \web@relation\leqslant{2A7D}%
    \web@relation\geqslant{2A7E}%
}{}

% --- alphabets and styles ---------------------------------------------------

% tex4ht maps \mathcal to mathvariant="script", which MathJax draws in a
% rounder script face; ask for the TeX calligraphic face instead. (For
% letters only)
\protected\def\mathcal#1{{\HCode{<mi mathvariant="-tex-calligraphic">#1</mi>}}}

% \bm comes out like \mathbf (bold upright); mark it for style.bold_symbols
\@ifpackageloaded{bm}{%
    \let\web@bm\bm
    \protected\def\bm#1{{\HCode{<mstyle class="bm">}\web@bm{#1}\HCode{</mstyle>}}}%
}{}

% tex4ht drops \textstyle (\textstyle\sum_i in a display); mark where it
% starts, and style.textstyle sets the rest of its row in text style.
% (\protected also because amsmath's \genfrac expands its style argument.
% Not \displaystyle, which tex4ht's own alignments and \genfrac use)
\let\web@textstyle\textstyle
\protected\def\textstyle{\web@textstyle\HCode{<mi class="textstyle"></mi>}}

% cases' table: tex4ht marks it only when mathtools is loaded (tables.array
% sets its spacing)
\@ifpackageloaded{amsmath}{%
    \ConfigureEnv{cases}{\HCode{<mrow class="cases">}}{\HCode{</mrow>}}{}{}%
}{}

% --- scripts ----------------------------------------------------------------

% tex4ht's _ and ^ take a single token, so x_\mathrm{op} and
% \underbrace{..}_\text{..} lose the argument. Make each font command a
% brace group, and have _ and ^ expand a font command (or an alias of one,
% such as \newcommand\mr{\mathrm}) before taking their token. (These stay
% expandable, for that)
\newcommand\web@bracefont[1]{%
    \ifdefined#1%
        \expandafter\let\csname web@orig\string#1\endcsname#1%
        \edef#1##1{{\expandafter\noexpand\csname web@orig\string#1\endcsname{##1}}}%
    \fi
}
\web@bracefont\mathrm
\web@bracefont\mathbf
\web@bracefont\mathsf
\web@bracefont\mathit
\web@bracefont\mathtt
\web@bracefont\mathbb
\web@bracefont\mathfrak
\web@bracefont\text
\ExplSyntaxOn
\clist_const:Nn \c__web_script_fonts_clist
    { \mathrm, \mathbf, \mathcal, \mathfrak, \mathsf, \mathbb, \mathit, \mathtt, \text, \bm }
\tl_new:N \l__web_script_target_tl
\bool_new:N \l__web_script_font_bool
% the command an alias stands for: while #1 is a macro without parameters
% whose replacement is a single token, that token
\cs_new_protected:Npn \__web_resolve_alias:N #1
  {
    \tl_set:Nn \l__web_script_target_tl {#1}
    \bool_lazy_and:nnT
      { \token_if_macro_p:N #1 }
      { \str_if_eq_p:ee { \cs_parameter_spec:N #1 } { } }
      {
        \exp_args:No \tl_if_single_token:nT {#1}
          { \exp_last_unbraced:No \__web_resolve_alias:N {#1} }
      }
  }
\cs_new_protected:Npn \__web_if_script_font:N #1
  {
    \bool_set_false:N \l__web_script_font_bool
    \clist_map_inline:Nn \c__web_script_fonts_clist
      {
        \token_if_eq_meaning:NNT #1 ##1
          { \bool_set_true:N \l__web_script_font_bool \clist_map_break: }
      }
  }
% #1: what _ or ^ does next, with the token
\cs_new_protected:Npn \__web_script_fix:n #1
  {
    \cs_set:Npn \__web_script:w {#1}
    \peek_after:Nw \__web_script_check:
  }
% only a macro can be a font command or an alias of one, and then it's safe
% to take it as an argument
\cs_new_protected:Npn \__web_script_check:
  { \token_if_macro:NTF \l_peek_token { \__web_script_check:N } { \__web_script:w } }
\cs_new_protected:Npn \__web_script_check:N #1
  {
    \__web_resolve_alias:N #1
    \exp_last_unbraced:NV \__web_if_script_font:N \l__web_script_target_tl
    \bool_if:NTF \l__web_script_font_bool
      { \exp_after:wN \__web_script:w \tex_romannumeral:D - `0 #1 }
      { \__web_script:w #1 }
  }
% x_.. goes through tex4ht's \S:b; x^.._.. through \SUP:SUB
\cs_set_eq:Nc \__web_tex_sub: { S:b }
\cs_set:cpn { S:b } { \__web_script_fix:n \__web_tex_sub: }
\cs_set_eq:Nc \__web_tex_supsub:n { SUP:SUB }
\cs_set:cpn { SUP:SUB } #1 { \__web_script_fix:n { \__web_tex_supsub:n #1 } }
\ExplSyntaxOff

% --- accents ----------------------------------------------------------------

% tex4ht writes \widehat and \hat alike; mark the wide accents, which stretch
% across their base (accents.widths keeps the others at their own width, as
% in TeX)
\Configure{accent}\widehat\widehat{{}{}}
    {}
    {\HCode{<mover accent="true" class="wide"><mrow>}#2\HCode{</mrow><mo accent="true">\string^</mo></mover>}}
\Configure{accent}\widetilde\widetilde{{}{}}
    {}
    {\HCode{<mover accent="true" class="wide"><mrow>}#2\HCode{</mrow><mo accent="true">\string~</mo></mover>}}

\makeatother
